\begin{terminologyList}
    \term[Адаптивное уточнение]{процесс локального увеличения детализации разреженной сетки в тех областях, где текущая аппроксимация оказывается недостаточно точной}
    \term[Базисная функция]{локальная функция, связанная с узлом сетки и определяющая вклад этого узла в итоговый интерполянт}
    \term[Векторное поле]{отображение, сопоставляющее каждой точке области вектор, например скорость изменения состояния динамической системы}
    \term[Входная точка сетки (Entry Point)]{начальный узел или стартовая конфигурация уровней и индексов, от которой начинается локальное построение адаптивной ветви разреженной сетки}
    \term[Дифференциальное уравнение]{уравнение, связывающее неизвестную функцию и её производные; в работе используется для задания динамики системы через правую часть}
    \term[Иерархический коэффициент (коэффициент избытка)]{разность между точным значением функции в узле и значением, восстановленным по уже построенной части сетки}
    \term[Интерполянт]{приближённое представление функции, позволяющее вычислять её значения по заранее построенной системе узлов и коэффициентов}
    \term[Квадратичный базис]{вариант базисной функции, в котором локальный вклад зависит квадратично от расстояния до центра узла; используется в работе как основной режим аппроксимации}
    \term[Контрольная точка (Anchor Point)]{заранее заданная точка области, в которой дополнительно проверяется ошибка интерполяции для принятия решения о продолжении адаптивного уточнения}
    \term[Параллельный алгоритм]{алгоритм, в котором независимые части вычислений выполняются одновременно на нескольких вычислительных потоках или ядрах процессора}
    \term[Разреженная сетка (Sparse Grid)]{иерархическая система узлов, предназначенная для аппроксимации многомерных функций с меньшим числом точек по сравнению с полной тензорной сеткой}
    \term[Рекурсивный обход]{способ обработки иерархической структуры данных, при котором вычисление в текущем узле дополняется вычислением во всех релевантных дочерних узлах}
    \term[Тензорная сетка]{полная многомерная сетка, получаемая декартовым произведением одномерных сеток; характеризуется быстрым ростом числа узлов с увеличением размерности}
    \term[Численное интегрирование]{приближённое вычисление решения дифференциального уравнения на конечном интервале времени с помощью дискретной схемы}
    \term[gRPC]{фреймворк удалённого вызова процедур, использующий protobuf-контракты и HTTP/2 для обмена структурированными сообщениями между клиентом и сервером}
    \term[Go]{компилируемый язык программирования, использованный в работе для реализации вычислительного ядра и серверного слоя}
    \term[Flutter]{кроссплатформенный UI-фреймворк, использованный для создания клиентского приложения визуализации}
\end{terminologyList}